% Generated draft: requires mathematical and source review.
\documentclass[11pt]{article}
\usepackage[margin=25mm]{geometry}
\usepackage{fontspec}
\setmainfont{Latin Modern Roman}
\usepackage{amsmath,amssymb,mathtools,unicode-math}
\setmathfont{Latin Modern Math}
\usepackage{enumitem,needspace,xurl,hyperref}
\setlength{\parindent}{0pt}
\setlength{\parskip}{5pt}
\setlength{\emergencystretch}{3em}
\newcommand{\N}{\mathbb{N}}
\newcommand{\Z}{\mathbb{Z}}
\newcommand{\Q}{\mathbb{Q}}
\newcommand{\R}{\mathbb{R}}
\newcommand{\C}{\mathbb{C}}
\newcommand{\F}{\mathbb{F}}
\newcommand{\A}{\mathbb{A}}
\newcommand{\PP}{\mathbb{P}}
\newcommand{\E}{\mathbb{E}}
\newcommand{\eps}{\varepsilon}
\newcommand{\dd}{\,\mathrm{d}}
\begin{document}
\section*{UnsolvedMath ID: 20000861}
\textbf{Source problem code:} catalogue entry
\subsection{Definitions and mathematical statement}
This record is Chapter D, ``A saliency map in primary visual cortex'', a Li Zhaoping (University College London) talk summary and discussion in AIM notes dated October 24, 2003. It is a contextual section rather than a determinate mathematical open problem. The source describes a recurrent intracortical V1 network transforming contrast input through contextual connections; firing rates signal saliency despite feature tuning, with the strongest response at a location supplying the score. Symbolically one may write $S(x)=\max_j r_j(x)$, without inferring a dynamical equation from that readout.

The discussion's qualifications are substantive. The speaker says all V1 cells contribute, other visual areas have no role in this model, and the theory concerns bottom-up saliency under an idealization disabling top-down feedback, for example shortly after exposure or under anesthesia. Responses are discussed after one time constant, not under a specified steady-state theorem. Long-range excitation is anisotropic, though no connection kernel is supplied. Color and motion are not included, natural-image experiments were not performed because of computational limits, and the endstopping question has no recorded answer. Steve questions extending visual-search saliency to three-dimensional surface pop-out; the speaker says the model does not explain those surface properties. Dana proposes an SNR explanation of the same data and questions blocked experimental designs; this is a participant's concern, not a verified finding that the data are false. The speaker says the same principle applies across scales without giving a multiscale rule, and says disynaptic inhibitory delay is needed to remove symmetry-breaking behavior. These are historical transcript statements, not independently verified physiological conclusions.

The aid adds a separate reduced-circuit theorem for commuting symmetric effective excitation and disynaptic-inhibition operators. Such operators admit a common orthogonal modal decomposition. The summary says each effective eigenvalue pair yields an exact scalar delay margin and onset frequency, but it gives neither the scalar delay equation nor the formulas or admissible eigenvalue assumptions. Hence no numerical margin is reconstructed. Conditional on the stated model's stability hypotheses and modal thresholds $\tau_j\in(0,\infty]$, the reported network classification is: if every $\tau_j=\infty$, exponential stability holds for each finite common delay; otherwise $\tau_* =\min_j\tau_j$ determines the first critical modal subspace. ``For each finite delay'' does not assert a uniform decay constant over unbounded delay.

Where the circuit additionally defines an invariant homogeneous mode and a nonuniform complement, the aid classifies that first critical space as nonuniform only if all minimizing modes are nonuniform, homogeneous if all minimizers are homogeneous, and mixed if both occur. Symmetric commutation alone does not identify the constant vector as an eigenmode, so the homogeneous-mode structure is an additional necessary convention. This is a conditional linear critical-space interpretation, not a proved nonlinear symmetry-breaking bifurcation or a validation of the transcript's delay claim.

The second reported lemma concerns finite rank-only observations. If V1-response scores and SNR scores on finitely many stimuli induce the same preorder, including ties, their distinct score values are related by a strictly increasing relabeling. With an unspecified monotone score-to-behavior link, such rank-only data cannot distinguish them. The lemma neither specifies an SNR formula nor excludes discrimination using calibrated amplitudes, timing or an independently fixed link. The source supplies no full network equations, delay law, dataset or behavioral observation model.
\subsection{Short English statement}
This is a 2003 discussion of a restricted bottom-up V1 saliency model, not one precise open problem. The aid adds a conditional modal-delay classification and a rank-identifiability observation. Their missing circuit equations and empirical links prevent treating them as a validation of the transcript's physiological claims.
\subsection{Sources}
\begin{itemize}
\item Local UnsolvedMath registry, permanent ID $20000861$, source code \texttt{catalogue entry}, supplied complete chapter transcript and research summary; canonical local record \texttt{aim-biology-notes.json}, \texttt{notes[25]}.
\item American Institute of Mathematics, \emph{Inference and Prediction in Neocortical Circuits}, Chapter D, ``A saliency map in primary visual cortex'', Li Zhaoping, University College London. The supplied verification identifies the notes as an AIM web-page hard copy dated October 24, 2003: \url{https://aimath.org/WWN/brain/brain.pdf}. No PDF was retrieved.
\item Discussion attributions retained as supplied: Bill, Steve and Dana; the transcript's damaged ``Tries-mann data'' reference has no full bibliographic identification. The original literature field is empty. No full authoritative citation for the auxiliary delay or rank lemmas is supplied.
\item Research attempt 1: context only; problem not a problem at the run, while registry status is open and classification \texttt{OPEN-TRIAGE}; supplied rigor label ``proved'', novelty confidence ``low''. The machine-generated aid is not peer reviewed and requires independent verification before citation.
\end{itemize}
\end{document}
